% Source: https://artofproblemsolving.com/community/c3h2087754_eisirrational_and_dchenmathcounts_2019_mock_mathcounts
% Stable identity: aops-topic-2087754-q01
For odd \(n\), \(\gcd(n,8)=1\), so \(\dfrac{1}{n}\) is purely periodic in base \(8\), and its minimal period is the least \(k\ge 1\) with \(n\mid 8^k-1\), that is, the order of \(8\) modulo \(n\). We need this order to be exactly \(4\).

The order divides \(4\) exactly when \(n\mid 8^4-1=4095=3^2\cdot 5\cdot 7\cdot 13\). The order divides \(2\) (period \(1\) or \(2\)) exactly when \(n\mid 8^2-1=63=3^2\cdot 7\). So the valid \(n\) are the divisors of \(4095\) that do not divide \(63\). (Every divisor of \(63\) also divides \(4095\).)

The student's error is twofold: the condition is on \(n\), not on the numerator \(x\), and the divisors of \(63\) (such as \(1,3,7,9,21,63\)) give periods \(1\) or \(2\), not \(4\), while \(4095\) itself does work (\(\dfrac{1}{4095}=0.\overline{0001}_8\)).

Using \(\sigma\) for the sum of divisors,
\[
\sigma(4095)=\sigma(9)\sigma(5)\sigma(7)\sigma(13)=13\cdot 6\cdot 8\cdot 14=8736,\qquad \sigma(63)=13\cdot 8=104.
\]
All divisors of \(4095\) are odd, so the required sum is
\[
8736-104=\boxed{8632}.
\]
